%   Macros on size
\pgfmathsetmacro{\cardwidth}{18}
\pgfmathsetmacro{\cardheight}{25.5}
\pgfmathsetmacro{\imagewidth}{\cardwidth}
\pgfmathsetmacro{\imageheight}{0.75*\cardheight}

% Macros for text
\newcommand{\EndRound}{\textbf{At end of round}, }
\newcommand{\imageSize}{1.5cm}
\newcommand{\imageInSize}{1.25cm}
\newcommand{\CostR}{\includegraphics[height = \imageSize]{CostR} }
\newcommand{\CostB}{\includegraphics[height = \imageSize]{CostB} }
\newcommand{\CostY}{\includegraphics[height = \imageSize]{CostY} }
\newcommand{\CostW}{\includegraphics[height = \imageSize]{CostW} }
\newcommand{\Desert}{\includegraphics[height = \imageSize]{Desert} }
\newcommand{\Marsh}{\includegraphics[height = \imageSize]{Marsh} }
\newcommand{\Ocean}{\includegraphics[height = \imageSize]{Ocean} }
\newcommand{\Jungle}{\includegraphics[height = \imageSize]{Jungle} }
\newcommand{\Population}{\includegraphics[height = \imageSize]{population} }
\newcommand{\Trait}{\includegraphics[height = \imageSize]{trait} }
\newcommand{\Gene}{\includegraphics[height = \imageInSize]{gene} }
\newcommand{\Disaster}{\includegraphics[height = \imageSize]{disaster} }
\newcommand{\Epidemic}{\includegraphics[height = \imageSize]{epidemic} }
\newcommand{\StartRound}{\includegraphics[height = \imageSize] {startRound} }
\newcommand{\WhenGain}{\includegraphics[height = \imageSize] {gain} }
\newcommand{\OnceRound}{\includegraphics[height = \imageSize] {onceRound}: }
\newcommand{\Feed}{\includegraphics[height = \imageSize] {Gene} }
\newcommand{\Another}{\includegraphics[height = \imageSize]{another} }
\newcommand{\vp}{\includegraphics[height = \imageInSize] {vp} }

%   Defines the card shape
\def\shapeCard{(0,0) rectangle (\cardwidth,\cardheight)}

%   Add Rounded corners.
\tikzset{
    cardcorners/.style={
        rounded corners=0.2cm
    }
}

% Draws the border and snips everything to be within the card area.
\newcommand{\cardborder}{
    %\clip[cardcorners] \shapeCard;
    %\draw[cardcorners] \shapeCard;
}

% Adds the stats for the card. The stats input in order:
% Cost, Strength, Prosperity, Loyalty
\newcommand{\cardcost}[3]{
		%\node[left] at (\cardwidth*.20,\cardheight*.8) {#1\CostR};
		\ifthenelse{ #1> 0}{\node[left] at (\cardwidth*.20,\cardheight*.8) {\Huge{#1}\CostR};}{}
		\ifthenelse{ #2> 0}{\node[left] at (\cardwidth*.20,\cardheight*.7) {\Huge#2\CostB};}{}
		\ifthenelse{ #3> 0}{\node[left] at (\cardwidth*.20,\cardheight*.6) {\Huge#3\CostY};}{}
}

\newcommand{\cardcatacolor}[2]{
	\node[right] at (\cardwidth*.7,\cardheight*.85) {\includegraphics[height = .7cm]{#1}};
	\node[right] at (\cardwidth*.7,\cardheight*.75) {\includegraphics[height = .7cm]{#2}};
}

\newcommand{\cardbiome}[1]{
	\node[left] at (\cardwidth*.50,\cardheight*.1) {\includegraphics[height = .7cm]{#1}};
}

\newcommand{\cardcata}[2]{
	\ifthenelse{ #1> 0}{\node[left] at (\cardwidth*.20,\cardheight*.8) {\Huge#1\Disaster};}{}
	\ifthenelse{ #2> 0}{\node[left] at (\cardwidth*.20,\cardheight*.7) {\Huge#2\Epidemic};}{}
}

\newcommand{\cardextra}[2] {
	\ifthenelse{ #1> 0}{\node[right] at (\cardwidth*.80,\cardheight*.8) {#1\Population};}{}
	\ifthenelse{ #2> 0}{\node[right] at (\cardwidth*.80,\cardheight*.7) {#2\Trait};}{}
	
	
}

% Adds the name, image and description to the card.
% The input is name, image_name
\newcommand{\cardinfo}[2]{
		% Add the title with some shadow
		\node[text width=14.5cm, below right] at (\cardwidth*.2,\cardheight*.95) {
       \shadowtext{\parbox{8.5cm}{\begin{spacing}{3.6}\textbf{\hspace{-1.5cm}\uppercase{\Huge #1}}\end{spacing}}}
    };
		% Add the image to the card.
		\node at (\cardwidth*.5, \cardheight*.7) {\includegraphics[height = 6.75cm, width = 9cm, keepaspectratio]{#2}};
}

% Adds the description and boxes it.
\newcommand{\cardcontent}[1]{
	\node[draw, rectangle, ultra thick, text width=15.5cm, below] at (\cardwidth*.5, 0.5*\cardheight) {\textrm{\Huge#1}};
}
%%%%%%%%%%%%%%%%%%%%%%%%%%%%%%%%%%
% Legacy commands

\newcommand{\cardbackground}[1]{
		\clip[cardcorners] \shapeCard;
		%\draw[fill=white] (\cardwidth-.2,\cardheight-.2)--(.2,\cardheight-.2)--(.2,.2)--(\cardwidth-.2,.2)--cycle;
    %\node at (\cardwidth/2, \cardheight/2) {\includegraphics[height = 9.5cm]{#1}};
		\node[rotate=270] at (5,6.15) {\textbf{\large{\uppercase{#1}}}};
		\node[rotate=0] at (5,8) {\includegraphics[height = .95cm]{#1}};
}



\newcommand{\cardhp}[1]{
		\node[left] at (\cardwidth-\stripwidth*.45,\cardheight-5.4*\stripwidth) {\textbf{\LARGE #1} \Health};
}


\newcommand{\cardimage}[1]{
	\node at (\cardwidth*.4, \cardheight*.7) {\includegraphics[height = 2.5cm, width = 3cm, keepaspectratio]{#1}};
}

\newcommand{\cardvp}[1]{
		\node at (\cardwidth-.7,\cardheight-.5) {\textbf{\LARGE #1\vp}};
}

\newcommand{\first}[1]{
	\node at (.2*\cardwidth, .45*\cardheight) {\includegraphics[height=1cm]{#1}};
}
\newcommand{\second}[1]{
	\node at (.4*\cardwidth, .45*\cardheight) {\includegraphics[height=1cm]{#1}};
}
\newcommand{\third}[1]{
	\node at (.6*\cardwidth, .45*\cardheight) {\includegraphics[height=1cm]{#1}};
}
\newcommand{\fourth}[1]{
	\node at (.8*\cardwidth, .45*\cardheight) {\includegraphics[height=1cm]{#1}};
}

\newcommand{\general}[1]{
\ifthenelse{\equal{#1}{1}\OR \equal{#1}{2}}{\draw[thick] (-1,.8) -- (2.5, .8) -- (2.5, -1); \node at (1.25, .27) {\includegraphics[height=1cm]{Player}};}{}
\ifthenelse{\equal{#1}{2}}{\draw[thick] (3.5,-1) -- (3.5, .8) -- (7, .8);\node at (4.75, .27) {\includegraphics[height=1cm]{Player}};}{}
}
